\scanpage{103}
\mspage{\handwritten{6}}
obvious sense by a finite number, nor even by a finite number of infinit\unsure{e} series of maps such as $k_{\deleted{\unsure{1}}\handwritten{\ell}}^i$, $\mathrm{inv}_i$, $\overset{i\deleted{\unsure{\ell}}}{*}_\ell$, the compatibility arrows in the pentagon, and the like. More precisely still, I am convinced that none of the functors $\underline{C}_n\to\underline{C}_{n+1}$ is an equivalence, which amounts to saying that the structures of increasing order form a\deleted{\unsure{n}} strictly increasing sequence - at every step, there is actual extra structure added.%
% slanted
\margin{\handwritten{\uline{\unsure{NB}} That $\overset{i}{*}_\ell$ is \unsure{of} order $\ell$ and not $\unsure{i|\ell-1}$ is heuristically \unsure{clear}, but will \unsure{need} a proof none the less!}}
 This is perhaps evident beforehand to topologists in the know, but I confess that for the time being it isn't to me, in terms uniquely of the somewhat formal description I will make of the passage of $\underline{C}_n$ to $\underline{C}_{n+1}$. \deleted{which}\inserted{\handwritten{This step}} theoretically is all that remains to be done, in order to achieve an explicit construction of the \deleted{notion of an} structure species of an \deleted{\unsure{$\infty$}}$\infty$-category (besides the definition of standart fiber products) - without having to get involved still less lost in the technical intricacies of ever messier diagrams to\deleted{\unsure{,}} write down, with increasing order of the structure\deleted{\unsure{s}} to be added \ldots

\handwritten{6.}\quad In the outline of a method of construction for the structure species, there has not been any explicit mention so far \unsure{of|af} the topological motivation behind the whole approach, which could wrongly give the impression of being a purely algebraic one. However, topological considerations alone are giving me the clue both for the description of the so-called standart fiber products, and of the inductive step allowing to wind up from $\underline{C}_n$ to $\underline{C}_{n+1}$ \unsure{?|.} The heuristics \deleted{\unsure{th}}\handwritten{indeed} of the present approach is simple enough, and suggested by the starting task, to define pertinent structure on the \inserted{system of} sets $F_i(X)$ of ``homotopies'' of arbitrary order, associated to an arbitrary topological space. In effect, the functors \deleted{\unsure{$X\longrightarrow F$}}
\[
X\longmapsto F_i(X)
\]
are representable by spaces $D_i$, which are easily seen to be $i$-cells. The source and target maps $s^i_1,t^i_1:F_i(X)\rightrightarrows F_{i-1}(X)$ are transposed to maps, which I may denote by the same \deleted{\unsure{set}}\handwritten{lett}ers
\[
s^i_1,t^i_1:\quad D_{i-1}\rightrightarrows D_i\,.
\]
Handling around a little, one easily convinces oneself that all the main structural items on $F_*(X)$ which one is figuring out in succession, such as the degeneracy maps $k_1^i$, the inversion maps $\mathrm{inv}_i$, the composition $v\bullet u=v*_1u$ for $i$-objects, etc, are all transposed of similar maps which are defined between the cells\deleted{\unsure{s}} $D_i$, or which go from such cells to certain spaces, deduced from these by gluing them together - the most evident example \inserted{in this respect} being the composition of paths, which is transposed of a map from the unit segment $\inserted{I=}D_1$ into $(I,1)\amalg_e(I,0)$, having preassigned values on the \unsure{endpoints|andpoints} of $I$ \handwritten{(}which correspond in fact to the images of the two maps $s_1^1$, $\unsure{t_1 1|t_1^1}:D_0={}\text{one point}\rightrightarrows D_1=I$\handwritten{)}.
% archive sheet number: 102
